\printedpage{108}
% head: A. Grothendieck : Dérivateurs, Ch. XII, Édition des Universités de Montpellier II et Paris VII
(\footnote[107]{careful, this is not a commutative square. Say rather that $(\mathcal{X}/x_0)_u \simeq \int \begin{tikzcd}[ampersand replacement=\&, row sep=0.4em, column sep=0.4em, baseline=(\tikzcdmatrixname-2-1.base)] \& {} \\ {} \ar[ur]\ar[dr] \& \\ \& {} \end{tikzcd}$ is smooth over $\Phi$ with $W$-aspheric fibers (if $f_0$, $f_1$, $f_2$ are smooth), hence $W$-aspheric since $\Phi$ is.}), we must express that $(\mathcal{X}/x_0)_u$ is $W$-aspheric, knowing that $(X_1/x_1)_u$ and $(X_2/x_2)_u$ are. So this is equivalent to % Resolved from PDF leaf 31: the original diagram has exactly these two unlabelled arrows and no vertex labels.
\[
i_1(x_0,u) \in W \qquad (\text{or } i_2(x_0,u) \in W),
\]
hence also (assuming W(7 bis)) to
\msp{page 157}
\[
\alpha_1(x_0,u) \in W, \qquad \text{i.e. } (X_0/x_0)_u \ W\text{-aspheric.}
\]
To say that this is true for every $u$ in $Y/y$, means that $f_0 : X_0 \longrightarrow Y$ is smooth at $u$ \texttt{[rather at $x_0$]}.

So we have proved a), in the more precise form:

\textbf{Corollary 6.} \textit{Let $x \in \Ob \mathcal{X}$ (where $\mathcal{X} = \int \begin{tikzcd}[ampersand replacement=\&, row sep=tiny, column sep=small, baseline=(\tikzcdmatrixname-2-1.base)] \& X_1 \\ X_0 \ar[ur]\ar[dr] \& \\ \& X_2 \end{tikzcd}$.) We want to express the $W$-smoothness of}
\[
\varphi_Q : \mathcal{X} \longrightarrow Y
\]
\textit{at $x$.}

\begin{enumerate}[label={}, leftmargin=2.5em]
\item[\textit{a)}] \textit{If $x \in \Ob X_i$ ($i \in \{1,2\}$), $\varphi_Q$ is $W$-smooth at $x$ if and only if $f_i : X_i \longrightarrow Y$ is.}
\item[\textit{b)}] \textit{Assume} \textup{W(7 bis)} (\footnote[108]{Unnecessary for the forward direction, assuming \texttt{[incomplete sentence]}.})\textit{. Let $x_0 \in \Ob X_0$, and suppose $f_1$, $f_2$ smooth at $x_1$, $x_2$. Then $\varphi_Q$ is $W$-smooth at $x_0$ if and only if $f_0 : X_0 \longrightarrow Y$ is.}
\end{enumerate}

\textit{So} (assuming W(7 bis)) \textit{in all three cases $x \in \Ob X_i$ ($i \in \{0,1,2\}$) we find that $\varphi_Q$ is $W$-smooth at $x$ if and only if $f_i : X_i \longrightarrow Y$ is} \texttt{[not really for $i=0$]}.

\cercle{b} As for b) of theorem 4, it is lemma 2 above. So everything is proved!

\msp{page 158}

\subsection*{3.5\quad Applications to amalgamated sums in a category $X^{\wedge}$ (with $X \in \Ob \Cat$).}

First a characterization of amalgamated sums involving an open immersion:

\textbf{Proposition 15.} \textit{Let}
% foot: Version 7 décembre 2002
\printedpage{109}
% head: A. Grothendieck : Dérivateurs, Ch. XII, Édition des Universités de Montpellier II et Paris VII
\[
Q : \qquad
\begin{tikzcd}
U \ar[r, "i"] \ar[d, "p"'] & X \ar[d, "q"] \\
U' \ar[r, "i'"] & X'
\end{tikzcd}
\]
\textit{be a commutative square in} $\Cat$\textit{. For this square to be cocartesian in} $\Cat$\textit{, with $i$ (hence $i'$) an open immersion, it is necessary and sufficient that the following conditions be satisfied:}

\begin{enumerate}[label={}, leftmargin=2.5em]
\item[\textit{a)}] \textit{$Q$ is cartesian.}
\item[\textit{b)}] \textit{$i'$ is an open immersion (hence $i$ too).}
\item[\textit{c)}] \textit{If $T' = X' \backslash U'$, and $T = X \times_{X'} T' = X \backslash U$, $T \longrightarrow T'$ is an isomorphism.}
\item[\textit{d)}] \textit{For every $t \in T$, if $t'$ is its image in $T'$, the functor induced by $p$}
\[
p/t : \underbrace{U/t}_{= V_t} \longrightarrow \underbrace{U'/t'}_{= V'_{t'}}
\]
\textit{satisfies $(p/t)_!(e_{U/t}) \simeq e_{U'/t'}$, or what amounts to the same, the ``colocal fibers'' of $p/t$, $v' \backslash V_t$, with $v'$ in $V'_{t'}$, are} $0$\textit{-connected.}
\end{enumerate}

\msp{page 159}
Immediate proof - this should rather go back to §1.

\textbf{Corollary 1.} \textit{Let $Q$ be a square that is} cartesian (cf. page 158), \textit{with $i'$ an} open immersion \textit{(hence $i$ too), and $q$} $W$-smooth \textit{(hence $p$ too). We assume that $W$ satisfies the axiom} W(5) \textit{of connectedness} (\footnote[109]{So it suffices $q$ $W_0$-smooth, where $W_0 = \{f \in \Fl \Cat \mid \pi_0(f) \text{ bijective}\}$.})\textit{. Let $T' = X' \backslash U'$, $T = X \backslash U$, suppose that $T' \longrightarrow T$ is an isomorphism. Then $Q$ is cocartesian.}

It suffices to prove condition d) of prop. 15. Let us consider the square
\[
\begin{tikzcd}
U/t \ar[r, hook, "i/t"] \ar[d, "p/t"'] & X/t \ar[d, "q/t"] \\
U'/t' \ar[r, hook, "i'/t'"] & X'/t',
\end{tikzcd}
\tag*{$(*)$}
\]
I say that it is \textit{cartesian.} This follows from the sole assumption that $Q$ is cartesian, without any smoothness assumption on $q$ nor that $i'$ is an open immersion, and can be seen e.g. by looking at the cubical diagram:
% foot: Version 7 décembre 2002
\printedpage{110}
% head: A. Grothendieck : Dérivateurs, Ch. XII, Édition des Universités de Montpellier II et Paris VII
\[
\begin{tikzcd}[row sep=small, column sep=small]
& U/t \ar[rr] \ar[dl] \ar[dd] & & X/t \ar[dl] \ar[dd] \\
U'/t' \ar[rr, crossing over] \ar[dd] & & X'/t' & \\
& U \ar[rr] \ar[dl] & & X \ar[dl] \\
U' \ar[rr] & & X' \ar[from=uu, crossing over] &
\end{tikzcd}
\]
(\footnote[110]{NB Here $t$ is an arbitrary element of $\Ob X$, $t'$ its image in $X'$.}). In this diagram, the front and back face of the cube is cartesian, by definition of $U/t$ and of $U'/t'$
\msp{page 160}
respectively, and the base is by assumption. It follows that the top face of the cube is also a cartesian square.

This seen, since $q$ is smooth, the same is true of $q/t$ in the cartesian square $(*)$, hence also of $p/t$ by base change. On the other hand, the smoothness of $q$ implies that $q/t$ has $W$-aspheric fibers, hence by base change the same is true of $p/t$ (which is thus a universal $W$-equivalence, i.e. a trivial $W$-fibration). Moreover, because of the smoothness of $p/t$, for $v'$ in $U'/t'$, the inclusion % [sic] commme -> comme, repeated letter
\[
(U/t)_{v'} \longrightarrow v' \backslash (U/t)
\]
is in $W$, hence (by W(5)) defines a bijection on the $\pi_0$. Since $(U/t)_{v'}$ is $W$-aspheric by what precedes, it follows (by W(5)) that its $\pi_0$ has cardinality 1, hence the same is true of $v' \backslash (U/t)$, i.e. the latter is 0-connected. Now that is condition d) of proposition 15, which was to be established.

\msp{page 161}
\textbf{Corollary 2.} \textit{Let $Q$ be a cartesian square} (p. 158), \textit{with $i'$ (hence also $i$) an open immersion, and $q$ (hence also $p$) a} local isomorphism. \textit{Let $T = X \backslash U$, $T' = X' \backslash U'$ (hence $T \simeq X \times_{X'} T'$), and suppose $T' \longrightarrow T$ an isomorphism. Then $Q$ is cocartesian.}

This is a special case of corollary 1, but with no explicit reference to a localizer $W$. We can take $W = W_{\infty}$, or $W = W_0$, to reduce to corollary 1. But it is more reasonable to say that in the cartesian square $(*)$, the assumption that $q$ is a local isomorphism implies that $q/t$ is an isomorphism, hence $p/t$ is too. It does indeed follow that the $v' \backslash (U/t)$ are 0-connected, this category being isomorphic to $v' \backslash (U'/t')$, hence has an initial object. % [sic] pendre -> prendre, missing letter

Now let $X$ be an object of Cat, and let us consider a cocartesian diagram in $X^{\wedge}$:
\[
\begin{tikzcd}
G \ar[r, "i"] \ar[d, "p"'] & F \ar[d, "q"] \\
G' \ar[r, "i'"] & F',
\end{tikzcd}
\tag*{$(*)$}
\]
% foot: Version 7 décembre 2002
\printedpage{111}
% head: A. Grothendieck : Dérivateurs, Ch. XII, Édition des Universités de Montpellier II et Paris VII
\msp{page 162}
where $i$ is a monomorphism. We know that the cocartesian character of the square is checked ``argument by argument'' by the cocartesian character of
\[
\begin{tikzcd}
G(x) \ar[r, "i(x)"] \ar[d, "p(x)"'] & F(x) \ar[d, "q(x)"] \\
G'(x) \ar[r, "i'(x)"] & F'(x),
\end{tikzcd}
\tag*{$(*)$}
\]
in Ens, $\forall\, x$ in $X$. Now such a cocartesian square in Ens, with $i(x)$ \texttt{[a]} monomorphism, is also cartesian. This implies that the square $Q$ in $X^{\wedge}$ is also cartesian.

Let us consider, on the other hand, the canonical functor
\[
\varphi_X : F \longmapsto X/F : X^{\wedge} \longrightarrow \mathrm{Cat}.
\]
It is better, moreover, to look at it as a functor
\[
\varphi_X : X^{\wedge} \longrightarrow \mathrm{Cat}/X,
\]
and I recall, for the record, that this functor is fully faithful, and that its essential image consists of the categories $X'$ over $X$ such that
\[
X' \longrightarrow X
\]
\msp{page 163}
is a local isomorphism, or what amounts to the same, that $X'$ is \textit{fibered} over $X$ with discrete fibers:
\[
\varphi_X : X^{\wedge} \xrightarrow{\ \sim\ } \underbrace{\Fibdisc\, X}_{\substack{\text{fibered categories}\\ \text{with discrete fibers}}} \ .
\]
Let us note that every arrow $X' \longrightarrow X''$ in $\Fibdisc\, X$ is a local isomorphism (tautological on $\begin{tikzcd}[ampersand replacement=\&, row sep=small, column sep=tiny, baseline=(\tikzcdmatrixname-1-1.base)] X'/x' \ar[rr] \ar[dr, "\text{iso.}"'] \& \& X''/x'' \ar[dl, "\text{iso.}"] \\ \& X/x \& \end{tikzcd}$), which shows that if we have a diagram
\[
\begin{tikzcd}[row sep=small, column sep=small]
X_1 \ar[dr] & & X_2 \ar[dl] \\
& X_0 &
\end{tikzcd}
\]
in $\Fibdisc\, X$, the fiber product $X_3 = X_1 \times_{X_0} X_2$ in Cat is such that $X_3 \longrightarrow X_2$ is a local isomorphism, hence also $X_3 \longrightarrow X$. That is to say $\Fibdisc(X)$ is stable under fiber products. This implies that fiber products exist in $\Fibdisc\, X$, and that they are the fiber products in Cat, or in Cat$/X$.
% foot: Version 7 décembre 2002

\printedpage{112}
% head: A. Grothendieck : Dérivateurs, Ch. XII, Édition des Universités de Montpellier II et Paris VII
\msp{page 164}
Likewise for finite products in $\Fibdisc\, X$, they are the finite products in Cat$/X$. In short

\textbf{Lemma 1.} \textit{The full subcategory} $\Fibdisc\, X$ \textit{of} $\Cat/X$ \textit{is stable under finite} $\varprojlim$\textit{, hence the} $\varprojlim$ \textit{exist in} $\Fibdisc\, X$ \textit{and the inclusion functor}
\[
\Fibdisc\, X \longrightarrow \Cat/X
\]
\textit{commutes with them, and in particular transforms cartesian squares of} $\Fibdisc\, X$ \textit{into cartesian squares of} $\Cat/X$\textit{, or what amounts to the same, of} $\Cat$\textit{. Moreover, by virtue of the equivalence}
\[
\varphi_X : X^{\wedge} \xrightarrow{\ \sim\ } \Fibdisc\, X \hookrightarrow \Cat/X,
\]
\textit{the functor $\varphi_X$ from $X^{\wedge}$ to} $\Cat/X$ \textit{is left exact, in particular it transforms cartesian squares of $X^{\wedge}$ into cartesian squares of} $\Cat/X$\textit{, or again of} $\Cat$\textit{.} % [sic] exacte -> exact, agreement with "le foncteur"

\msp{page 165}
Let us also note that being left exact, the functor $\varphi_X$ from $X^{\wedge}$ to Cat$/X$ transforms monomorphisms into monomorphisms. On the other hand:

\textbf{Lemma 2.} \textit{The monomorphisms which are local isomorphisms are exactly the open immersions. So in} $\Fibdisc\, X$\textit{, the monomorphisms are the open immersions, and the functor $\varphi_X : F \longmapsto X/F$ transforms monomorphisms into open immersions.}

Let us then return to the cocartesian diagram $(*)$ (p. 161), and to the corresponding diagram in Cat (and even in $\Fibdisc\, X$)
\[
\begin{tikzcd}
X/G \ar[r, hook, "X/i"] \ar[d, "X/p"'] & X/F \ar[d, "X/q"] & T \ar[l, hook', dotted] \ar[d, dotted] \\
X/G' \ar[r, hook, "X/i'"] & X/F' & T'. \ar[l, hook', dotted]
\end{tikzcd}
\]
Since $(*)$ is cartesian, the preceding square is cartesian, by what precedes. I want to prove that it is also cocartesian
\msp{page 166}
(in Cat$/X$, or in Cat, it amounts to the same), and for this I apply corollary 2 of proposition 15 (p. 161), taking into account that $X/q$ is a local isomorphism, and $X/i'$ (being a local isomorphism \texttt{[and a]} monomorphism) an open immersion. It suffices to prove that, $T'$ and $T$ being the closed subcategories complementary to the opens $X/G'$ of $X/F'$ and to $X/G$ of $X/F$, that $T \longrightarrow T'$ is \texttt{[an]} isomorphism.

{\small \texttt{[Now, $F$ and $G$, and $F'$ and $G'$, but not $p$ and $q$, exchange their roles.]}} % [sic] leur rôles -> leurs rôles, agreement

The verification is immediate: $T$ is the category of arrows $x \longrightarrow G$ in $X^{\wedge}$ (with source $x$ in $X$) which do not factor through $G$ \texttt{[rather through $F$]}, $T'$ that of arrows $x \longrightarrow G'$
% foot: Version 7 décembre 2002